\section{Escalado del cómputo en el momento de la inferencia (Test-Time Compute Scaling)}

\subsection{Idea central}

Ya se ha demostrado que el training-time scaling (aumentar los datos de entrenamiento y el cómputo) es eficaz. Test-Time Compute Scaling aplica la misma idea a la etapa de inferencia: \textbf{invertir más cómputo en el momento de la inferencia a cambio de una salida de mayor calidad}.

\subsection{Majority Voting (votación mayoritaria)}

El método más simple de test-time scaling:

\begin{enumerate}
    \item Para la misma pregunta (por ejemplo, un problema de matemáticas), se le pide al modelo que genere $N$ respuestas distintas (usando una temperatura distinta de cero para obtener diversidad)
    \item Se toma la respuesta que aparece con mayor frecuencia como salida final
\end{enumerate}

Intuición: debido a la aleatoriedad que introduce el muestreo top-p, el modelo ocasionalmente elige un token equivocado que hace que toda la cadena de razonamiento se desvíe. Al muestrear varias veces, es muy probable que la respuesta correcta aparezca con mayor frecuencia.

\subsection{Best-of-N (selección óptima)}

Una versión mejorada de Majority Voting:

\begin{enumerate}
    \item Generar $N$ respuestas distintas
    \item Usar un \textbf{Reward Model} o un \textbf{Judge LLM} para puntuar cada respuesta
    \item Elegir la respuesta con la puntuación más alta
\end{enumerate}

En comparación con Majority Voting, Best-of-N aprovecha una señal de evaluación externa, por lo que suele funcionar mejor, aunque requiere ejecutar un modelo de evaluación adicional.

\subsection{Process Reward Model (PRM)}

Un método más refinado: en lugar de evaluar la respuesta completa, evalúa \textbf{cada paso del razonamiento}:

\begin{importantbox}{Flujo de trabajo del Process Reward Model}
\begin{enumerate}
    \item El LLM genera una solución parcial (partial solution)
    \item El PRM puntúa cada paso: la puntuación representa "la probabilidad de que ese paso finalmente conduzca a la respuesta correcta"
    \item Se elige el paso con la puntuación más alta para seguir expandiéndolo
    \item Se repite el proceso anterior para construir la solución completa paso a paso
\end{enumerate}
Esto da lugar a un mecanismo de \textbf{búsqueda en árbol} (tree search): se busca en el espacio de los pasos de razonamiento, y el PRM guía la dirección de la búsqueda.
\end{importantbox}

El profesor demostró este proceso con un ejemplo de multiplicación en base 8:

\begin{itemize}
    \item Attempt 1: el primer paso es correcto, pero el segundo se calcula en base decimal en lugar de en base 8, y el PRM lo rechaza
    \item Attempt 2: tiene un paso correcto más, pero al final comete el mismo error, y el PRM lo rechaza de nuevo
    \item Attempt 4: todos los pasos son correctos, y el PRM lo aprueba
\end{itemize}

\subsection{Resultados experimentales de Scaling}

Los experimentos del equipo de Hugging Face muestran que:

\begin{itemize}
    \item Usando Llama 3.2 1B y 3B (modelos pequeños), al aumentar el número de generaciones por pregunta
    \item en un benchmark de matemáticas, un modelo pequeño + test-time scaling puede \textbf{superar} a Llama 3.1 8B e incluso a Llama 3.1 70B
\end{itemize}

\begin{importantbox}{El equilibrio entre la velocidad de inferencia y la calidad de la salida}
En esencia, Test-Time Compute Scaling intercambia \textbf{tiempo de inferencia} (inference speed) por \textbf{calidad de la salida} (output quality). Esto significa que un modelo pequeño y eficiente junto con un presupuesto de inferencia suficiente puede igualar o incluso superar a modelos de un tamaño mucho mayor en tareas específicas. Esto tiene una importancia considerable para los escenarios con recursos limitados.
\end{importantbox}

\begin{knowledgebox}{De dónde viene el Reward Model}
El Reward Model que se usa en Test-Time Compute Scaling puede reutilizar el reward model entrenado en la etapa de Preference Alignment (dentro del proceso de entrenamiento con PPO). La entrada de este modelo es texto, y su salida es una puntuación de calidad entre 0 y 1. También es posible sustituir el reward model entrenado específicamente por un Judge LLM (mediante prompting).
\end{knowledgebox}

\subsection{Resumen de la sección}

Test-Time Compute Scaling es una de las tendencias más importantes de principios de 2025; su idea central es invertir más cómputo en el momento de la inferencia para mejorar la calidad de la salida. Desde el simple Majority Voting, pasando por Best-of-N, hasta la búsqueda a nivel de paso basada en el Process Reward Model, la complejidad del método y su efectividad aumentan progresivamente. Los experimentos demuestran que un modelo pequeño + test-time scaling puede superar a modelos más grandes.

\newpage
